\documentclass[aspectratio=169,11pt]{beamer}

% =========================================================
% PACKAGES
% =========================================================

\usepackage[T1]{fontenc}
\usepackage{lmodern}

\usepackage{graphicx}
\usepackage{amsmath}
\usepackage{amssymb}
\usepackage{tikz}
\usepackage{xcolor}

\usepackage{styles/couleurs}

\graphicspath{{images/}}

% =========================================================
% STYLE GÉNÉRAL
% =========================================================

\setbeamertemplate{navigation symbols}{}

\setbeamertemplate{footline}{%
  \hfill
  \scriptsize\insertframenumber
  \hspace{0.4cm}
  \vspace{0.15cm}
}

\setbeamercolor{normal text}{
  fg=black,
  bg=white
}

\setbeamercolor{frametitle}{
  fg=TitreC,
  bg=white
}

\setbeamerfont{frametitle}{
  series=\bfseries,
  size=\Large
}

\setbeamersize{
  text margin left=0.5cm,
  text margin right=0.5cm
}

% =========================================================
% OUTILS
% =========================================================

% ---------------------------------------------------------
% Image encadrée
% Aucun trim, aucun clip :
% l'image entière est toujours conservée.
% ---------------------------------------------------------

\newcommand{\imagecadree}[2][]{%
  \begingroup
  \setlength{\fboxsep}{1pt}%
  \setlength{\fboxrule}{0.4pt}%
  \fcolorbox{black!25}{white}{%
    \includegraphics[#1]{#2}%
  }%
  \endgroup
}

% ---------------------------------------------------------
% Correction éclair : géométrie
% Format volontairement compact pour laisser
% le maximum de place à la production.
% ---------------------------------------------------------

\newcommand{\correctiongeo}{%
\begin{beamercolorbox}[
  wd=\textwidth,
  sep=0.12cm,
  rounded=true
]{block body}

\footnotesize

\textbf{Correction éclair :}
\(\widehat{ACB}=40^\circ\),
\(\widehat{BCD}=60^\circ\),
\(40^\circ+60^\circ+80^\circ=180^\circ\).

Donc \(A,C,E\) sont alignés et
\(AE=4+2{,}4=6{,}4\) cm.

\end{beamercolorbox}

\vspace{0.08cm}
}

% ---------------------------------------------------------
% Correction éclair : allumettes
% ---------------------------------------------------------

\newcommand{\correctionallumettes}{%
\begin{beamercolorbox}[
  wd=\textwidth,
  sep=0.12cm,
  rounded=true
]{block body}

\footnotesize

\textbf{Correction éclair :}
si \(n\) est le numéro de l'étape,

\[
5+4(n-1)=4n+1.
\]

\end{beamercolorbox}

\vspace{0.05cm}
}

% ---------------------------------------------------------
% Question orale
% ---------------------------------------------------------

\newcommand{\questionorale}{%
\begin{center}
  {\color{ActionC}\bfseries
  Qu'est-ce que vous remarquez ?}
\end{center}

\vspace{-0.15cm}
}

% =========================================================
% DOCUMENT
% =========================================================

\begin{document}

% =========================================================
% DIAPO 1
% TITRE
% =========================================================

\begin{frame}[plain]

\vfill

\begin{center}

{\Huge\bfseries\color{TitreC}
J'écris des maths}

\vspace{0.3cm}

{\Large\color{SavoirC}
Étude de productions d'élèves}

\vspace{1cm}

{\large
\textbf{\(+\)}
quelque chose qui m'aide
\qquad
\textbf{\(?\)}
quelque chose qui m'interroge
\qquad
\textbf{\(\rightarrow\)}
quelque chose que je modifierais
}

\end{center}

\vfill

\end{frame}


% =========================================================
% DIAPO 2
% ÉNONCÉ GÉOMÉTRIE
% =========================================================

\begin{frame}[t]
\frametitle{Problème de géométrie}

\vspace{-0.15cm}

\begin{columns}[T,onlytextwidth]

% ---------------------------------------------------------
% COLONNE ÉNONCÉ
% ---------------------------------------------------------

\begin{column}{0.53\textwidth}

\small

\(ABC\) est un triangle rectangle en \(B\) tel que

\[
\widehat{BAC}=50^\circ
\qquad\text{et}\qquad
AC=4\text{ cm}.
\]

\(BCD\) est un triangle équilatéral.

\medskip

\(CDE\) est un triangle tel que

\[
\widehat{DCE}=80^\circ
\qquad\text{et}\qquad
CE=2{,}4\text{ cm}.
\]

\medskip

{\bfseries Calculer la longueur \(AE\).}

\end{column}

% ---------------------------------------------------------
% COLONNE FIGURE
% ---------------------------------------------------------

\begin{column}{0.44\textwidth}

\centering

\imagecadree[
  width=\linewidth,
  height=0.52\textheight,
  keepaspectratio
]{problème de longueur figure.png}

\end{column}

\end{columns}

\vspace{0.25cm}

\begin{beamercolorbox}[
  wd=\textwidth,
  sep=0.18cm,
  rounded=true
]{block body}

\small
\textbf{Correction éclair}

\[
\widehat{ACB}=40^\circ,
\qquad
\widehat{BCD}=60^\circ,
\qquad
40^\circ+60^\circ+80^\circ=180^\circ.
\]

Donc \(A,C,E\) sont alignés et

\[
AE=4+2{,}4=6{,}4\text{ cm}.
\]

\end{beamercolorbox}

\end{frame}


% =========================================================
% DIAPO 3
% PRODUCTION 1
% =========================================================

\begin{frame}[t]
\frametitle{Production 1}

\correctiongeo
\questionorale

\vfill

\begin{center}

\imagecadree[
  width=0.96\textwidth,
  height=0.59\textheight,
  keepaspectratio
]{production 1.png}

\end{center}

\vfill

\end{frame}


% =========================================================
% DIAPO 4
% PRODUCTION 3 - PARTIE 1
% =========================================================

\begin{frame}[t]
\frametitle{Production 3}

\correctiongeo
\questionorale

\vfill

\begin{center}

\imagecadree[
  width=0.94\textwidth,
  height=0.64\textheight,
  keepaspectratio
]{production 3 (1).png}

\end{center}

\vfill

\end{frame}


% =========================================================
% DIAPO 5
% PRODUCTION 3 - PARTIE 2
% =========================================================

\begin{frame}[t]
\frametitle{Production 3 -- suite}

\correctiongeo
\questionorale

\vfill

\begin{center}

\imagecadree[
  width=0.94\textwidth,
  height=0.64\textheight,
  keepaspectratio
]{production 3 (2).png}

\end{center}

\vfill

\end{frame}


% =========================================================
% DIAPO 6
% PRODUCTION 4
% =========================================================

\begin{frame}[t]
\frametitle{Production 4 -- brouillon}

\correctiongeo
\questionorale

\vfill

\begin{center}

\imagecadree[
  width=0.94\textwidth,
  height=0.65\textheight,
  keepaspectratio
]{production 4.png}

\end{center}

\vfill

\end{frame}


% =========================================================
% DIAPO 7
% PRODUCTION 5 - PARTIE 1
% =========================================================

\begin{frame}[t]
\frametitle{Production 5}

\correctiongeo
\questionorale

\vfill

\begin{center}

\imagecadree[
  width=0.95\textwidth,
  height=0.65\textheight,
  keepaspectratio
]{production 5 (1).png}

\end{center}

\vfill

\end{frame}


% =========================================================
% DIAPO 8
% PRODUCTION 5 - PARTIE 2
% =========================================================

\begin{frame}[t]
\frametitle{Production 5 -- suite}

\correctiongeo
\questionorale

\vfill

\begin{center}

\imagecadree[
  width=0.96\textwidth,
  height=0.65\textheight,
  keepaspectratio
]{production 5 (2).png}

\end{center}

\vfill

\end{frame}


% =========================================================
% DIAPO 9
% PRODUCTION 6
% =========================================================

\begin{frame}[t]
\frametitle{Production 6}

\correctiongeo
\questionorale

\vfill

\begin{center}

% Production très verticale :
% on privilégie ici la hauteur.
\imagecadree[
  width=0.75\textwidth,
  height=0.66\textheight,
  keepaspectratio
]{production 6.png}

\end{center}

\vfill

\end{frame}


% =========================================================
% DIAPO 10
% ÉNONCÉ ALLUMETTES
% =========================================================

\begin{frame}[t]
\frametitle{Problème des allumettes}

\vspace{0.1cm}

\begin{columns}[T,onlytextwidth]

% ---------------------------------------------------------
% IMAGE
% ---------------------------------------------------------

\begin{column}{0.49\textwidth}

\centering

\imagecadree[
  width=\linewidth,
  height=0.49\textheight,
  keepaspectratio
]{allumettes maisons.png}

\end{column}

% ---------------------------------------------------------
% ÉNONCÉ
% ---------------------------------------------------------

\begin{column}{0.48\textwidth}

\small

On construit des maisons avec des allumettes
comme sur la figure.

\medskip

Il fallait :

\begin{itemize}
  \item trouver le nombre d'allumettes aux étapes
  \(4\), \(10\) et \(2007\) ;
  \item donner une méthode valable pour une
  \textbf{étape quelconque}.
\end{itemize}

\end{column}

\end{columns}

\vspace{0.25cm}

\begin{beamercolorbox}[
  wd=\textwidth,
  sep=0.18cm,
  rounded=true
]{block body}

\small
\textbf{Correction éclair}

\[
5+4(n-1)=4n+1
\]

\[
\text{Étape 4 : }17
\qquad
\text{Étape 10 : }41
\qquad
\text{Étape 2007 : }8029
\]

\end{beamercolorbox}

\end{frame}


% =========================================================
% DIAPO 11
% RÉPONSE 2
% =========================================================

\begin{frame}[t]
\frametitle{Réponse 2}

\correctionallumettes
\questionorale

\vfill

\begin{center}

\imagecadree[
  width=0.94\textwidth,
  height=0.54\textheight,
  keepaspectratio
]{reponse 2.png}

\end{center}

\vfill

\end{frame}


% =========================================================
% DIAPO 12
% RÉPONSE 4
% =========================================================

\begin{frame}[t]
\frametitle{Réponse 4}

\correctionallumettes
\questionorale

\vfill

\begin{center}

\imagecadree[
  width=0.94\textwidth,
  height=0.54\textheight,
  keepaspectratio
]{reponse 4.png}

\end{center}

\vfill

\end{frame}


% =========================================================
% DIAPO 13
% RÉPONSE 6
% =========================================================

\begin{frame}[t]
\frametitle{Réponse 6}

\correctionallumettes
\questionorale

\vfill

\begin{center}

\imagecadree[
  width=0.90\textwidth,
  height=0.57\textheight,
  keepaspectratio
]{reponse 6.png}

\end{center}

\vfill

\end{frame}


% =========================================================
% DIAPO 14
% RÉPONSE 7
% =========================================================

\begin{frame}[t]
\frametitle{Réponse 7}

\correctionallumettes
\questionorale

\vfill

\begin{center}

\imagecadree[
  width=0.91\textwidth,
  height=0.59\textheight,
  keepaspectratio
]{reponse 7.png}

\end{center}

\vfill

\end{frame}


% =========================================================
% DIAPO 15
% RÉPONSE 8
% =========================================================

\begin{frame}[t]
\frametitle{Réponse 8}

\correctionallumettes
\questionorale

\vfill

\begin{center}

\imagecadree[
  width=0.82\textwidth,
  height=0.59\textheight,
  keepaspectratio
]{reponse 8.png}

\end{center}

\vfill

\end{frame}

\end{document}